This study compares interpretable and nonlinear classifiers for identifying existing gallstone status from demographic, laboratory, and bioimpedance measurements. A secondary analysis of 319 records with 38 candidate predictors evaluated logistic regression, random forest, XGBoost, LightGBM, and stacking using leave-one-patient-out (LOPO) internal validation, implemented as record-level leave-one-out under a one-record-per-person assumption. Imputation, recursive CART-based selection of 20 predictors, and transformation were fitted within training folds. Each model produced one held-out prediction per record. Logistic regression achieved the highest pooled ROC--AUC (0.874), whereas random forest achieved the highest accuracy (0.799), sensitivity (0.839), and F1 (0.808) among predictive models. Stacking achieved the highest specificity (0.772). These results distinguish discrimination from classification at the fixed threshold of 0.5 and support logistic regression as an interpretable research baseline. Recurrent errors across model families indicate limitations beyond aggregate scores. A small single-source cohort, overlapping training folds, retrospective model comparison, and absent calibration and external validation restrict inference. The findings concern existing disease status and do not establish prospective prediction or clinical utility.
